\begin{textsample}{POS Dim 1 – rs – Score 50.00 – rs/rs\_em\_52.txt}  \label{ex:f1_pos_204}
3.3.3 Caracterização das Ciências Humanas Sociais e Aplicadas Para a área das CHS – composta por Ensino Religioso, Filosofia, Geografia, História e Sociologia – a BNCC propõe ampliar e aprofundar as aprendizagens do Ensino Fundamental, voltadas para uma formação geral pautada pela ética como compromisso educativo que observa \textbf{ideais} e práticas de justiça, solidariedade, autonomia, liberdade de pensamento e de escolha.
Essas constantes se organizam de modo a estudar, pesquisar, investigar, debater, refletir, problematizar e permitir produções, construções e proposições considerando -se as especificidades de cada território e região, bem como aspectos socioculturais, econômico-políticos e histórico-científicos dos diferentes povos.
Encontram -se expressas nas \textbf{categorias} \textbf{conceituais} de Tempo e Espaço ; Territórios e Fronteiras ; Indivíduo, Natureza, Sociedade, Cultura e Ética ; Política e Trabalho, assim estabelecidas pela BNCC : - A \textbf{categoria} de Tempo e Espaço visa garantir que os estudantes sejam capazes de, no Ensino Médio, analisar acontecimentos, comparar, observar e empreender continuidades e mudanças em diferentes sociedades, como prevê a BNCC, estudando suas instituições, as \textbf{razões} de suas desigualdades, de seus micro e macro conflitos e suas relações de poder.
A \textbf{categoria}, no entanto, não se esgota nessas dimensões, pois considera também a necessidade da superação de sentidos e de práticas \textbf{tanto} por parte da escola e dos currículos quanto das redes, dos sistemas e dos sujeitos envolvidos nos processos de aprendizagens. - Território é uma \textbf{categoria} das CHS que se refere, como indica a BNCC, a “ noções de lugar, região, fronteira e, especialmente, os limites políticos e administrativos de cidades ”, Estados e Nações ( BRASIL, 2018a ).
Além das dimensões concretas, o conceito de território \textbf{implica} o simbólico e as abstrações da instituição sociocultural e a concepção de “ poder, jurisdição, administração e soberania, dimensões que expressam a diversidade das relações sociais ”. - Outra \textbf{categoria} que compõe a conceitualidade das CHS, a Fronteira expressa, na proposta da BNCC ( BRASIL, 2018a ), a cultura, a \textbf{pluralidade} étnica, as identidades, as estruturas e organizações sociais particulares, grupais, bem como as características particulares das diversas manifestações.
Fronteira significa limites entre territórios e processos, propriedades e espaços, demarcação de espaços, tempos e \textbf{compreensões} físicas e \textbf{teóricas} ; também entre civilização e barbárie, alfabetizado e analfabeto, capaz e incapaz, branco e preto, capital e trabalho.
Fronteiras são geográficas e culturais, físicas e simbólicas e precisam ser pensadas nas suas implicações, estruturas, constituições e consequências. - Na \textbf{categoria} Indivíduo, Natureza, Sociedade, Cultura e Ética, as relações e os conceitos \textbf{complexos} que a envolvem apresentam um amplo campo de \textbf{compreensão}, \textbf{debate} e produção.
Aí se encontram, se relacionam e dialogam entre si a subjetividade, a objetividade e a intersubjetividade, a autocompreensão e a alteridade, o cuidado consigo e com o outro e com a vida, como elementos \textbf{primordiais} para a existência, a coexistência, a manutenção, defesa e continuidade das espécies e do ambiente.
Perguntas fundamentais são realizadas por pensadores da Filosofia, Sociologia, História, Geografia e Ensino Religioso, desde longos tempos até os dias ora compartilhados.
As CHS envolvem as problemáticas \textbf{inerentes} aos indivíduos em suas subjetividades e relações e se expressam na dimensão da Política e do Trabalho.
A política é \textbf{instância} de decisões, enquanto o trabalho é a ação de transformação da natureza e construção do mundo da vida.
Desse modo, as Ciências Humanas e Sociais Aplicadas no Ensino Médio têm como base “ as ideias de justiça, solidariedade, autonomia, liberdade de pensamento e de escolha, a \textbf{compreensão} e o reconhecimento das diferenças, o respeito aos direitos humanos e à interculturalidade e o combate aos preconceitos de qualquer natureza ” ( BRASIL, 2018a, p. 561 ), além de integrar conhecimentos que estão para além do eurocentrismo, do colonialismo e neo-colonialismo e, assim, indicam a necessidade de conhecimentos referentes à história dos povos negros e indígenas ( africanos, americanos, brasileiros e gaúchos ).
Nessa perspectiva, destacam -se as contribuições do pensamento descolonial e decolonial, nas suas diferentes \textbf{vertentes} \textbf{teóricas}, pois adquirem relevância não somente nos meios acadêmicos ao expandir -se por diferentes campos de conhecimento, \textbf{influenciando} distintas formas de atuação social, política e ética, mas, igualmente, no mundo da \textbf{ciência}, da economia e das fontes \textbf{constituintes} da história, da \textbf{filosofia}, da sociologia, das religiosidades e da geografia.
O conhecimento e o estudo da plêiade de pensadores e perspectivas \textbf{teóricas} vinculadas a essa linha de pensamento em suas \textbf{matrizes}, é \textbf{indispensável} para investigações, \textbf{análises} e produções na área de CHS e a constituição cidadã da sociedade brasileira e gaúcha.
Em consonância com as \textbf{competências} específicas da área de CHS elencadas na BNCC, tal \textbf{abordagem} propicia o entendimento de que o conhecimento, embora regulado por relações de poder e lógicas de dominação social, cultural, \textbf{econômica} e política na história da \textbf{humanidade}, pode e deve ser ressignificado a partir da \textbf{pluralidade} e diversidade na \textbf{interface} com as diversas concepções \textbf{teórico-metodológicas}, sobretudo, a partir de experiências coletivas e individuais vivenciadas em diferentes contextos, territórios e relações de poder.
Além disso, o desafio da \textbf{interdisciplinaridade} se coloca com urgência nas práticas docentes.
A escola \textbf{contemporânea} deve acompanhar as profundas transformações pelas quais passa a sociedade.
Neste momento histórico, em que a velocidade das informações e conhecimentos imprime o nosso agir social, professores e professoras devem ser capazes de inovar e \textbf{renovar} metas, objetivos e metodologias.
O trabalho pedagógico deve ajudar as \textbf{juventudes} na \textbf{sistematização} das muitas e variadas informações com as quais se confrontam, para que se transformem em conhecimento e favoreçam práticas protagonistas e formação de sujeitos \textbf{críticos}, participativos, capazes de pensar o mundo e de agir nele com argumentação e comunicação \textbf{dialógica}.
O mundo interconectado e \textbf{complexo} derruba a \textbf{tese} do conhecimento como ação isolada, singular e individual e afirma um mundo plural, capaz de debater ideias, trocar experiências e integrar os diferentes saberes presentes nos componentes curriculares.
A \textbf{interdisciplinaridade} coloca em novos níveis de comunicação, interação \textbf{reflexiva} e trabalho coletivo e pode resultar em \textbf{aulas} mais participativas, dinâmicas e significativas para a aprendizagem com encontro de diferentes saberes.
Educar parece ser também construir estratégias de relação que conectem os sujeitos aprendentes, estudantes e professores/as, porque os seres humanos nascem em suas relações e nelas se constituem ; fora delas, não existe vida humana.
Relação existe quando o outro aparece na relação não como o mesmo, como qualquer um, como comum, pois é da diferença que brota a possibilidade de relação e do reconhecimento que germina a qualidade da relação.
O professor é mediador das relações de aprendizagem e \textbf{colaborador} para construir condições para a \textbf{sistematização} das informações, para a orientação dos estudos, para a promoção do diálogo e da convivência que podem transformar os dados em conhecimento e conhecimento em práticas vitais, atitudes.
O trabalho educativo multi-inter-transdisciplinar compreende planejar conjuntamente, discutir coletivamente conteúdos, \textbf{métodos}, procedimentos e estratégias facilitadoras, atraentes e interconectadas na aquisição de conhecimentos.
Compreende, ainda, realização conjunta de atividades pedagógicas como \textbf{seminários}, rodas de conversa, bate-papos, \textbf{palestras} e ações de intervenção social com estudantes, investigações científicas, produções e comunicações dos resultados.
Aos professores que atuam nos componentes curriculares de Ensino Religioso, Filosofia, Geografia, História e Sociologia \textbf{sugere} -se o desenvolvimento ações pedagógicas integradoras e \textbf{articuladas}, que valorizem o protagonismo juvenil, busquem alcançar a preparação básica para a pesquisa científica, o desenvolvimento da cidadania e autonomia dos estudantes e \textbf{aprimorar} suas \textbf{capacidades} de \textbf{sistematizar} conhecimentos e apresentá -los na forma escrita e oral, de modo dialogado, bem como saber enfrentar os tensionamentos próprios das diversidades de \textbf{compreensões} e pensamentos.
Nas \textbf{seções} seguintes, são apresentadas temáticas e conceitos específicos dos componentes curriculares que integram as Ciências Humanas e Sociais Aplicadas.
Tais temáticas e conceitos aparecem apenas como sugestões – podendo ser adaptados, complementados, aperfeiçoados – de \textbf{caminhos} para o desenvolvimento das \textbf{competências} e habilidades da Área de Conhecimento, contribuindo, portanto, para a \textbf{compreensão} dos sentidos indicados pelo Ensino Médio e em incorporação pelo RCGEM.

% matched lemmas: abordagem, análise, aprimorar, articulado, aula, caminho, capacidade, categoria, ciência, colaborador, competência, complexo, compreensão, conceitual, constituinte, contemporâneo, crítico, debate, dialógico, econômico, filosofia, humanidade, ideal, implicar, indispensável, inerente, influenciar, instância, interdisciplinaridade, interface, juventude, matriz, método, palestra, pluralidade, primordial, razão, reflexivo, renovar, seminário, seção, sistematizar, sistematização, sugerir, tanto, tese, teórico, teórico-metodológico, vertente
\end{textsample}
